\documentclass[11pt,a4paper]{article}

% === 中文支持 ===
\usepackage{ctex}
\usepackage[T1]{fontenc}

% === 页面布局 ===
\usepackage[top=2.5cm, bottom=2.5cm, left=2.5cm, right=2.5cm]{geometry}

% === 表格 ===
\usepackage{booktabs}
\usepackage{array}
\usepackage{enumitem}
\usepackage{tabularx}

% === 数学符号 ===
\usepackage{amssymb}

% === 流程图（TikZ）===
\usepackage{tikz}
\usetikzlibrary{shapes.geometric, arrows.meta, positioning, calc, fit}
\usepackage{float}

\tikzset{
  arr/.style={thick, -Stealth},
  flowstep/.style={rectangle, rounded corners, draw=black, fill=white, thick,
    minimum width=2.6cm, minimum height=0.82cm, align=center, font=\small},
  decision/.style={diamond, draw=black, fill=white, thick, aspect=2.2,
    inner sep=1pt, minimum width=1.6cm, minimum height=1.6cm, align=center, font=\small},
  rolehuman/.style={ellipse, draw=blue!60, fill=blue!15, thick,
    minimum width=2.0cm, minimum height=0.72cm, align=center, font=\small},
  rolepm/.style={ellipse, draw=purple!60, fill=purple!15, thick,
    minimum width=2.0cm, minimum height=0.72cm, align=center, font=\small},
  roledev/.style={ellipse, draw=green!60, fill=green!15, thick,
    minimum width=2.0cm, minimum height=0.72cm, align=center, font=\small},
  rolesys/.style={ellipse, draw=gray!60, fill=gray!15, thick,
    minimum width=2.0cm, minimum height=0.72cm, align=center, font=\small},
  scanbox/.style={rectangle, rounded corners, draw=gray!70, fill=gray!10, thick,
    minimum width=1.9cm, minimum height=0.65cm, align=center, font=\scriptsize},
  statebox/.style={rectangle, rounded corners, thick, minimum width=1.85cm,
    minimum height=0.7cm, align=center, font=\scriptsize},
}

% === 颜色 ===
\usepackage{xcolor}
\definecolor{primary}{HTML}{1a1a2e}
\definecolor{accent}{HTML}{3b82f6}
\definecolor{success}{HTML}{22c55e}
\definecolor{danger}{HTML}{ef4444}
\definecolor{warning}{HTML}{f59e0b}
\definecolor{graytext}{HTML}{6b7280}
\definecolor{progress}{HTML}{8b5cf6}

% === 超链接 ===
\usepackage[hidelinks]{hyperref}
\hypersetup{colorlinks=true, linkcolor=primary, urlcolor=accent}

% === 标题样式 ===
\usepackage{titlesec}
\titleformat{\section}{\Large\bfseries\color{primary}}{}{0em}{}[\titlerule]
\titleformat{\subsection}{\large\bfseries\color{primary}}{}{0em}{}
\titleformat{\subsubsection}{\normalsize\bfseries\color{primary}}{}{0em}{}

% === 页眉页脚 ===
\usepackage{fancyhdr}
\pagestyle{fancy}
\fancyhf{}
\fancyhead[L]{\small\color{graytext} 企业级 AI Agent 应用商店 -- 总需求文档}
\fancyhead[R]{\small\color{graytext} MASTER v1.3}
\fancyfoot[C]{\small\color{graytext} \thepage}
\renewcommand{\headrulewidth}{0.4pt}

% === 代码 ===
\usepackage{listings}
\lstset{basicstyle=\ttfamily\footnotesize, breaklines=true, frame=single, showstringspaces=false}

% === 文档信息 ===
\title{
\vspace{-1em}
{\Huge\bfseries\color{primary} 企业级 AI Agent 应用商店}\\[0.8em]
{\Large\bfseries\color{accent} 总需求文档（MASTER）}\\[0.5em]
{\large 项目全部需求的完整汇总}
}
\author{
\textbf{PM AI}：万能产品小助手（OpenClaw）\\[0.3em]
\textbf{审阅}：魏子政
}
\date{2026-05-12 \quad 版本 v1.3}

\begin{document}

\maketitle
\thispagestyle{empty}

\begin{center}
\begin{tabular}{>{\bfseries}p{3cm} p{10cm}}
\toprule
文档性质 & 总需求文档（MASTER）-- 项目唯一的需求汇总 \\
更新规则 & 每次新增/变更需求后同步更新本文件 \\
与增量文档关系 & DEMAND-xxx 是每次需求的详细快照（冻结）；本文件是所有需求的完整汇总（活的） \\
技术栈 & Vue 3 + Element Plus + FastAPI + PostgreSQL + MinIO + OpenSearch + JWT + Docker Compose \\
\bottomrule
\end{tabular}
\end{center}

\vspace{1em}

\begin{abstract}
本文件是项目唯一的总需求文档，完整汇总所有需求的全貌。每个需求均包含描述、决策、任务分解、验收标准。

增量需求文档（DEMAND-xxx）记录每次需求的详细过程和决策（冻结快照），本文件汇总所有需求（活的 Backlog）。

每次生成 DEMAND-xxx 后，必须同步更新本文件。
\end{abstract}

\newpage
\tableofcontents
\newpage

% ============================================================
\section{需求总表}
% ============================================================

\begin{table}[h]
\centering
\small
\begin{tabular}{clp{4.5cm}cc}
\toprule
\textbf{编号} & \textbf{需求名称} & \textbf{摘要} & \textbf{阶段} & \textbf{状态} \\
\midrule
\multicolumn{5}{l}{\textbf{V0 先备阶段}} \\
\midrule
D-001 & 技术选型锁定 & 8 层技术栈，v1.0 精简为 v2.0 & V0 & \textcolor{success}{已完成} \\
D-002 & 智能体分工定义 & OpenClaw=需求+验收，Cursor=前后端+UI+环境+运维+测试 & V0 & \textcolor{success}{已完成} \\
D-003 & Skill/规则体系 & 两套智能体的上下文加载机制与文件清单 & V0 & \textcolor{success}{已完成} \\
D-004 & 文档体系与规范 & tex-pdf 格式、命名规范、版本管理、目录结构 & V0 & \textcolor{success}{已完成} \\
D-005 & 需求变更流程 & 5 步处理 + 3 级变更控制 + 4 级优先级 & V0 & \textcolor{success}{已完成} \\
D-006 & 自闭环机制 & 需求/开发/测试/文档/审计 5 个闭环 & V0 & \textcolor{success}{已完成} \\
D-007 & 开源方案调研 & Flow Nexus + Membrane，结论为自研 & V0 & \textcolor{success}{已完成} \\
\midrule
\multicolumn{5}{l}{\textbf{V1 正式需求阶段}} \\
\midrule
D-008 & 开发环境搭建 & Docker Compose + 前后端初始化 + 健康检查 & V1 & \textcolor{success}{已完成} \\
D-009 & 任务流转路径定义 & 魏子政 $\to$ OpenClaw $\to$ Cursor 的完整链路 & V1 & \textcolor{progress}{进行中} \\
D-010 & 多智能体协作方案 & OpenClaw/Claude/Codex/Cursor 的 skill 适配 & V1 & \textcolor{progress}{进行中} \\
D-011 & MVP 上传+审查可视化 & Skill zip 上传 + 三层扫描 + 审批工作台 & V1 & \textcolor{success}{已完成} \\
D-012 & 前端重新开发 & 正式 UI：市场/详情/我的应用/审批 & V1 & \textcolor{success}{已完成} \\
D-013 & 流转修复+安装能力 & 上传入库/审批工作台/zip下载/OpenClaw安装 & V1 & \textcolor{success}{已完成} \\
D-014 & 目录框架重构 & 清空旧代码+实习清单融合+目录蓝图 & V1 & \textcolor{success}{已完成} \\
D-015 & 邮箱认证注册 & 图形验证码；邮箱验证 MVP（响应体验证码，SMTP 待后续） & V1 & \textcolor{progress}{进行中} \\
D-016 & 前端UI品质升级 & 暗色主题、品牌色系、卡片式布局、响应式 & V1 & \textcolor{success}{已完成} \\
D-017 & 上架下架流程 & 下架/重新上架/重新提交完整流转（无独立 DEMAND 文件，见 D-015） & V1 & \textcolor{success}{已完成} \\
D-018 & 竞品分析 & Flow Nexus、Membrane、OpenAI GPT Store 等调研 & V1 & \textcolor{success}{已完成} \\
D-019 & SkillHub CLI & jsdhub：login/search/install/download（\texttt{pip install jsdhub}） & V1 & \textcolor{success}{已完成} \\
\bottomrule
\end{tabular}
\end{table}

% ============================================================
\section{V0 先备阶段需求（已完成）}
% ============================================================

\subsection{D-001：技术选型锁定}

\textbf{决策}：8 层技术栈锁定，v1.0 精简为 v2.0。
\begin{itemize}[nosep]
  \item 前端：Vue 3 + Element Plus + Vite + TypeScript + Pinia
  \item 后端：FastAPI + SQLAlchemy (async) + Pydantic
  \item 数据库：PostgreSQL 16（含 JSONB）
  \item 搜索：OpenSearch 2.x
  \item 存储：MinIO（S3 兼容）
  \item 认证：JWT（python-jose + passlib）
  \item 编排：Docker Compose
  \item 迁移：Alembic
\end{itemize}
\textbf{淘汰项}：Next.js/React、Spring/Java、Keycloak、Temporal、Kafka、MongoDB、K8s。
\textbf{对应文档}：DEMAND-000-v0-master, REF-002-project-master-plan。

\subsection{D-002：智能体分工定义}

\textbf{决策}：PM（OpenClaw）只做需求生成 + 拆解 + 独立验收。Cursor 全权负责前端（含 UI 设计）+ 后端 + 环境 + 运维 + 测试。
\textbf{对应文档}：DEMAND-001-role-redefine。

\subsection{D-003：Skill/规则体系搭建}

\textbf{交付物}：
\begin{itemize}[nosep]
  \item OpenClaw：4 个 SKILL.md（developer/pm/tester/ui-design）+ AGENTS.md + CLAUDE.md
  \item Cursor：1 个 CLAUDE.md + 7 个 .mdc（project-overview, backend-rules, frontend-rules, ui-design, api-conventions, testing, database）
\end{itemize}

\subsection{D-004：文档体系与规范}

\textbf{交付物}：tex-pdf 格式规范、命名规则、两层需求文档结构（MASTER + DEMAND-xxx）、目录分类（reference/ 仅放外部调研，demands/ 放需求+项目规范）。

\subsection{D-005：需求变更流程}

\textbf{流程}：需求新增 $\to$ PM 评估 $\to$ 拆解分配 $\to$ Cursor 执行 $\to$ PM 验收。3 级变更控制（P0/P1/P2），4 级优先级。

需求变更流程如下图所示：
\begin{figure}[H]
\centering
\begin{tikzpicture}[
    transform shape, scale=0.75,
    node distance=0.95cm and 0.55cm,
  ]
  \node[flowstep] (n1) {需求新增};
  \node[flowstep, right=1.2cm of n1] (n2) {PM 评估};
  \node[flowstep, right=1.2cm of n2] (n3) {拆解分配};
  \node[flowstep, right=1.2cm of n3] (n4) {Cursor 执行};
  \node[decision, right=1.0cm of n4] (n5) {PM\\验收};
  \node[flowstep, right=1.2cm of n5] (ok) {通过结案};

  \draw[arr] (n1) -- (n2);
  \draw[arr] (n2) to[bend left=18] node[above, sloped, font=\scriptsize] {P0 紧急} (n3);
  \draw[arr] (n2) -- node[above, font=\scriptsize] {P1 常规} (n3);
  \draw[arr] (n2) to[bend right=18] node[below, sloped, font=\scriptsize] {P2 优化} (n3);
  \draw[arr] (n3) -- (n4);
  \draw[arr] (n4) -- (n5);
  \draw[arr] (n5) -- node[above] {通过} (ok);
  \draw[arr] (n5.south) -- +(0,-0.35) -| node[pos=0.32, left, font=\scriptsize] {不通过} (n4.south);
\end{tikzpicture}
\caption{需求变更流程}
\label{fig:d005-change-flow}
\end{figure}

\subsection{D-006：自闭环机制}

5 个闭环：需求闭环、开发闭环、测试闭环、文档闭环、审计闭环。每个闭环有明确的触发条件、执行步骤、完成标准。

\subsection{D-007：开源方案调研}

\textbf{调研对象}：Flow Nexus（MCP 工具调用架构）+ Membrane（中间件代理）。
\textbf{结论}：两个方案均不完整，本项目需自研。可借鉴：分类体系、质量标准、模板部署、免密钥代理。

% ============================================================
\section{V1 正式需求阶段}
% ============================================================

\subsection{D-008：开发环境搭建与依赖安装}

\textbf{状态}：已完成。

\textbf{任务分解}：
\begin{enumerate}[nosep]
  \item 后端 FastAPI 项目初始化（main.py, config.py, database.py, /health 接口）
  \item 前端 Vue 3 + Vite + TS 项目初始化（Element Plus, Pinia, Vue Router）
  \item Docker Compose 编排（PostgreSQL 16 + MinIO + OpenSearch 2.x + Backend）
  \item 环境变量模板（.env.example）
\end{enumerate}

\textbf{验收结论}：
\begin{itemize}[nosep]
  \item[\checkmark] pip install -r requirements.txt 成功
  \item[\checkmark] cd frontend \&\& npm install 成功
  \item[\checkmark] uvicorn 启动成功，/health 返回 200
  \item[\checkmark] npm run build 构建成功（3.59s）
  \item[\checkmark] docker-compose.yml 四个服务配置完整
\end{itemize}

\textbf{对应文档}：DEMAND-002-dev-env-setup。

% ---------- D-009 ----------
\subsection{D-009：任务流转路径定义}

\textbf{状态}：进行中。
\textbf{提出者}：魏子政。
\textbf{优先级}：\textcolor{danger}{P0}。

\textbf{需求描述}：定义魏子政（甲方）$\to$ OpenClaw（PM AI）$\to$ Cursor（开发 AI）的完整任务流转路径（7 步流程）。核心新增：
\begin{itemize}[nosep]
  \item \textbf{拍板机制}：有歧义时 OpenClaw 必须先让魏子政拍板，再往下走
  \item \textbf{工具链锁定}：仅 OpenClaw + Cursor，Claude/Codex 暂不参与
\end{itemize}

\textbf{对应文档}：DEMAND-003-task-workflow（v0.2）。

任务流转路径如下图所示：
\begin{figure}[H]
\centering
\begin{tikzpicture}[
    transform shape,
    node distance=0.75cm and 0.9cm,
    scale=0.75,
    transform shape,
  ]
  \node[rolehuman] (a1) {魏子政\\[-2pt]\scriptsize （甲方）};
  \node[rolepm, right=of a1] (a2) {OpenClaw\\[-2pt]\scriptsize （PM AI）};
  \node[decision, right=of a2] (a3) {拍板\\确认};
  \node[roledev, right=of a3] (a4) {Cursor\\[-2pt]\scriptsize （开发 AI）};
  \node[rolesys, right=of a4] (a5) {代码审查};
  \node[rolepm, right=of a5] (a6) {PM 验收};
  \node[rolesys, right=of a6] (a7) {交付};

  \draw[arr] (a1) -- (a2);
  \draw[arr] (a2) -- (a3);
  \draw[arr] (a3) -- node[above, font=\scriptsize] {无歧义} (a4);
  \draw[arr] (a3.south) -- +(0,-0.95) -| node[pos=0.18, below, font=\scriptsize] {有歧义} (a1.south);
  \draw[arr] (a4) -- (a5);
  \draw[arr] (a5) -- (a6);
  \draw[arr] (a6) -- (a7);
\end{tikzpicture}
\caption{任务流转路径（7 步）}
\label{fig:d009-task-flow}
\end{figure}

% ---------- D-010 ----------
\subsection{D-010：多智能体协作方案}

\textbf{状态}：进行中。
\textbf{提出者}：魏子政。
\textbf{优先级}：\textcolor{warning}{P1}（当前仅 OpenClaw + Cursor，Claude/Codex 暂不参与）。

\textbf{需求描述}：设计多智能体 skill 适配方案。当前阶段工具链锁定为 OpenClaw + Cursor，Claude 和 Codex 作为备选，待后续需要时再接入。

\textbf{对应文档}：DEMAND-004-multi-agent-collaboration。

% ---------- D-011 ----------
\subsection{D-011：MVP 开发（上传 + 审查流程可视化）}

\textbf{状态}：已完成。
\textbf{提出者}：魏子政。
\textbf{优先级}：\textcolor{danger}{P0}。

\textbf{需求描述}：开发最小化 MVP，核心目标是跑通 Skill 上传和审查流程可视化。

\textbf{拍板决策}（魏子政已确认）：
\begin{itemize}[nosep]
  \item 三层安全扫描全上：Semgrep + ClamAV + qwen-flash（LLM 语义分析）
  \item 审查流程：管理员在审批工作台查看三层扫描结果（通过/不通过/告警详情），手动通过或驳回
  \item LLM 模型：通义千问 qwen-turbo（测试 API Key 已记录）
  \item API 地址：\texttt{https://dashscope.aliyuncs.com/compatible-mode/v1}
  \item 前端风格：黑白灰简约色系
  \item 初始数据：虚设 skill + 管理员账号（默认全部管理员权限）
  \item 交付物：局域网可访问地址 + 可查验的初始数据
\end{itemize}

\textbf{验收标准}：
\begin{itemize}[nosep]
  \item[\checkmark] 用户可上传 zip 包（Skill）
  \item[\checkmark] 上传后自动触发三层扫描（Semgrep + ClamAV + qwen-flash）
  \item[\checkmark] 管理员可在审批工作台查看扫描结果详情
  \item[\checkmark] 管理员可手动通过或驳回
  \item[\checkmark] 前端黑白灰简约色系
  \item[\checkmark] 局域网可访问，初始数据可查验
\end{itemize}

\textbf{对应文档}：DEMAND-005-mvp-upload-review。

% ---------- D-012 ----------
\subsection{D-012：前端重新开发}

\textbf{状态}：已完成。

\textbf{优先级}：\textcolor{danger}{P0}（承接 MVP 后的正式界面）。

\textbf{需求描述}：将 Demo 级前端重写为正式 UI，覆盖市场、详情、我的应用、审批工作台等核心页面。

\textbf{技术栈}：Vue 3 + Element Plus + vue-router + i18n。

\textbf{对应文档}：DEMAND-006-frontend-rewrite。

% ---------- D-013 ----------
\subsection{D-013：流转修复与安装能力}

\textbf{状态}：已完成。

\textbf{需求描述}：补齐上传入库、审批工作台操作闭环；支持 zip 下载；集成 OpenClaw 安装能力，形成可端到端验证的流转。

\textbf{对应文档}：DEMAND-007-flow-fix-and-install。

Skill 完整生命周期（上传、三层扫描、审批与上架）如下图所示：
\begin{figure}[H]
\centering
\begin{tikzpicture}[
    transform shape, scale=0.78,
    node distance=0.7cm and 0.65cm,
  ]
  \node[flowstep] (u) {用户上传 ZIP};
  \node[scanbox, below left=0.85cm and 0.18cm of u] (s1) {Semgrep};
  \node[scanbox, below=0.85cm of u] (s2) {ClamAV};
  \node[scanbox, below right=0.85cm and 0.18cm of u] (s3) {LLM 语义};
  \node[inner sep=3pt, draw=gray!50, dashed, rounded corners,
    fit=(s1)(s2)(s3), label={[font=\scriptsize,yshift=2pt]above:三层扫描}] (scanfit) {};
  \node[flowstep, below=0.72cm of scanfit] (agg) {扫描结果汇总};
  \node[flowstep, below=of agg] (adm) {管理员审批};
  \node[decision, below=of adm] (dec) {结果};
  \node[flowstep, below left=1.0cm and 0.45cm of dec] (rej) {驳回};
  \node[flowstep, below right=1.0cm and 0.45cm of dec] (pub) {上架发布};
  \node[flowstep, below=of pub] (use) {浏览 / 安装 / 下载};

  \draw[arr] (u) -- (s1);
  \draw[arr] (u) -- (s2);
  \draw[arr] (u) -- (s3);
  \draw[arr] (s1) |- (agg);
  \draw[arr] (s2) -- (agg);
  \draw[arr] (s3) |- (agg);
  \draw[arr] (agg) -- (adm);
  \draw[arr] (adm) -- (dec);
  \draw[arr] (dec.west) -- ++(-0.32,0) |- node[pos=0.46, left, font=\scriptsize] {驳回} (rej.north);
  \draw[arr] (dec.east) -- ++(0.32,0) |- node[pos=0.46, right, font=\scriptsize] {通过} (pub.north);
  \draw[arr] (rej.west) -- ++(-0.55,0) |- node[pos=0.62, left, font=\scriptsize] {重新提交} (u.west);
  \draw[arr] (pub) -- (use);
\end{tikzpicture}
\caption{Skill 完整生命周期流程}
\label{fig:skill-lifecycle}
\end{figure}

% ---------- D-014 ----------
\subsection{D-014：目录框架重构}

\textbf{状态}：已完成。

\textbf{需求描述}：清空旧代码噪声，融合实习清单与项目目录蓝图，统一仓库结构便于后续迭代。

\textbf{对应文档}：DEMAND-008-dir-framework-rebuild。

% ---------- D-015 ----------
\subsection{D-015：邮箱认证注册}

\textbf{状态}：进行中（MVP 模式：邮箱验证码在响应体返回，真实 SMTP 发送待后续）。

\textbf{优先级}：\textcolor{warning}{P1}。

\textbf{需求描述}：注册流程增加图形验证码；邮箱验证采用 MVP 实现（接口返回验证码），后续再接入 SMTP。

\textbf{对应文档}：DEMAND-009-email-auth-and-lifecycle。

% ---------- D-016 ----------
\subsection{D-016：前端 UI 品质升级}

\textbf{状态}：已完成。

\textbf{需求描述}：暗色主题、品牌色系、卡片式布局、响应式适配，使整体达到产品级应用商店观感。

\textbf{对应文档}：DEMAND-010-frontend-ui-upgrade。

% ---------- D-017 ----------
\subsection{D-017：上架与下架流程}

\textbf{状态}：已完成。

\textbf{需求描述}：支持下架、重新上架、重新提交等完整状态流转，与审批及版本策略衔接。

\textbf{对应文档}：无独立 DEMAND 文件（约定与实现见 DEMAND-009-email-auth-and-lifecycle 及相关代码变更）。

上架与下架状态流转如下图所示：
\begin{figure}[H]
\centering
\begin{tikzpicture}[
    transform shape, scale=0.78,
    node distance=1.05cm and 1.15cm,
  ]
  \node[statebox, fill=yellow!12, draw=warning!80!black] (st0) {草稿};
  \node[statebox, fill=blue!8, draw=blue!55, right=of st0] (st1) {已上传};
  \node[statebox, fill=progress!15, draw=progress!70, right=of st1] (st2) {扫描中};
  \node[statebox, fill=orange!12, draw=orange!70!black, right=of st2] (st3) {待审核};
  \node[statebox, fill=success!15, draw=success!70!black, below=1.5cm of st3] (st4) {已上架};
  \node[statebox, fill=gray!18, draw=gray!65, left=of st4] (st5) {已下架};

  \draw[arr] (st0) -- node[above, font=\scriptsize] {上传} (st1);
  \draw[arr] (st1) -- node[above, font=\scriptsize] {开始扫描} (st2);
  \draw[arr] (st2) -- node[above, font=\scriptsize] {扫描完成} (st3);
  \draw[arr] (st3) -- node[right, font=\scriptsize] {审批通过} (st4);
  \draw[arr] (st3.south) -- +(0,-0.58) -| node[pos=0.48, below, font=\scriptsize] {驳回} (st1.south);
  \draw[arr] (st4) -- node[above, font=\scriptsize] {手动下架} (st5);
  \draw[arr] (st5) -- node[below, font=\scriptsize] {重新上架} (st4);
\end{tikzpicture}
\caption{上架与下架状态机}
\label{fig:d017-state-machine}
\end{figure}

% ---------- D-018 ----------
\subsection{D-018：竞品分析}

\textbf{状态}：已完成（2026-05-11）。

\textbf{需求描述}：调研 Flow Nexus、Membrane、OpenAI GPT Store 等竞品能力与形态。

\textbf{结论摘要}：本项目需自研落地；可借鉴竞品在分类体系、质量标准、模板化部署等方面的设计。

\textbf{对应文档}：DEMAND-014-competitive-analysis。

% ---------- D-019 ----------
\subsection{D-019：SkillHub CLI（jsdhub）}

\textbf{状态}：已完成（2026-05-11）。

\textbf{需求描述}：提供命令行工具 \texttt{jsdhub}，支持 login、search、install、download 等操作，便于开发者从终端对接应用商店。

\textbf{交付形态}：Python 包名 \texttt{jsdhub}，安装方式为 \texttt{pip install jsdhub}。

\textbf{对应文档}：DEMAND-015-skillhub-cli。

% ============================================================
\section*{更新记录}
% ============================================================

\begin{tabular}{lll}
\toprule
版本 & 日期 & 变更内容 \\
\midrule
v1.3 & 2026-05-12 & 需求总表对齐 DEMAND 编号至 D-019；修正 D-016/D-017 条目；补充 D-012--D-019 章节；D-011 验收已完成；D-015 进行中说明 \\
v1.2 & 2026-05-06 & 新增 D-011 MVP 开发（上传+审查可视化），记录拍板决策 \\
v1.1 & 2026-05-06 & core/ 四份文档对齐 MASTER v1.0 修正（详见各文档更新记录）；新增文档一致性维护规则 \\
v1.0 & 2026-05-06 & 重构总需求文档：每个需求含完整描述/决策/任务/验收；新增 D-009/D-010 \\
v0.2 & 2026-05-06 & 新增 D-008 开发环境搭建（已完成） \\
v0.1 & 2026-05-06 & 初始版本 \\
\bottomrule
\end{tabular}

\end{document}
